\documentclass[titlepage,a4paper]{jsarticle}
\usepackage{../../../sty/import}% 各種パッケージインポート
\usepackage{../../../sty/title}% タイトルページの変更

%% タイトルページの変数
% レポートタイトル
\title{原子力材料と核燃料}
% 提出日
\expdate{\today}
% 科目名
\subject{原子力材料と核燃料}
% 分野
\class{情報経営システム工学分野}
% 学年
\grade{M1}
% 学籍番号
\mynumber{24336488}
% 記述者
\author{本間三暉}

\begin{document}
% titleページ作成
\maketitle
\section{金属の強化方法}
金属材料の強化は，転位の移動を妨げ，塑性変形を起こしにくくすることで行われる．
主な方法は以下のとおりである．

1．結晶粒微細化強化
結晶粒を小さくすると，転位の移動を妨げる結晶粒界が増加する．
そのため，転位が移動しにくくなり，材料の強度が高くなる．

2．固溶強化
母材に大きさの異なる元素を固溶させると，結晶格子にひずみが生じる．
このひずみ場が転位の移動を妨げるため，材料が強化される．
置換型固溶体と侵入型固溶体のいずれでも生じる．

3．析出強化
母相中に微細で硬い析出物を形成させる方法である．
析出物が転位の移動に対する障害となり，転位が析出物を切断または迂回するために大きな応力が必要となる．

4．加工硬化
圧延，引抜き，鍛造などの冷間加工によって塑性変形を加える方法である．
加工によって転位密度が増加し，転位同士が絡み合って移動しにくくなるため，強度および硬さが増加する．

以上のように，金属材料は結晶粒界，溶質原子，析出物および転位同士の相互作用を利用して強化される．

% 参考文献
\begin{thebibliography}{99}
\bibitem{} 原子力材料と核燃料 講義資料
\end{thebibliography}

\end{document}